% ============================================================
%  CHAPTER 8: TESTING AND EXPERIMENTAL METHODOLOGY
% ============================================================

\chapter{TESTING AND EXPERIMENTAL METHODOLOGY}
\label{chap:testing}

\section{Verification and Validation Strategy}
\label{sec:testing_strategy}

A comprehensive multi-tiered verification and validation (V\&V) strategy was executed to ensure that QuantumSafe VPN satisfies all functional specifications, cryptographic invariants, performance thresholds, and mobile operating system constraints. The testing methodology follows a structured V-Model progression spanning five discrete verification tiers:

\begin{enumerate}
    \item \textbf{Tier 1: Cryptographic Unit Testing}: Isolated mathematical verification of ML-KEM-1024, ML-DSA-65, Keccak-$f[1600]$ SHA3-256, and AES-256-GCM against known-answer test vectors.
    \item \textbf{Tier 2: Wire Protocol Conformance and Interoperability Testing}: Verification of binary serialization, length-prefixed framing, and cross-platform shared secret decapsulation between Android Kotlin and Linux Python runtimes.
    \item \textbf{Tier 3: Network Security and Adversary Injection Testing}: Active simulation of Man-in-the-Middle (MitM) attacks, ciphertext bit-flipping, packet replaying, and truncated framing.
    \item \textbf{Tier 4: System Integration and Tunneling Verification}: End-to-end evaluation of Android \texttt{VpnService} packet capture, socket loop protection via \texttt{protect()}, IP pool lease allocation, and Linux NAT masquerading.
    \item \textbf{Tier 5: Operational Usability and Robustness Testing}: Evaluation of UI state synchronization, background service persistence, abrupt network interface handovers, and DNS leak prevention.
\end{enumerate}

\section{Experimental Testbed Environment}
\label{sec:testbed_environment}

Testing was conducted across a dedicated physical hardware testbed configured to reflect production enterprise mobile conditions:

\begin{itemize}
    \item \textbf{Physical Mobile Client Node}:
    \begin{itemize}
        \item Device: Google Pixel / Android Reference Device.
        \item Operating System: Android 14 (API Level 34), non-rooted standard commercial image.
        \item Local Wi-Fi IP: \texttt{11.12.2.82} on IEEE 802.11ac 5\,GHz WLAN.
        \item Assigned Virtual VPN IP: \texttt{10.8.0.2/24}.
    \end{itemize}
    \item \textbf{Linux VPN Gateway Node}:
    \begin{itemize}
        \item Hardware: Dedicated x86\_64 host (Intel Core i7-12700H, 32\,GB DDR5 RAM).
        \item Operating System: Kali Linux 2024.2 (Kernel 6.6.9-amd64).
        \item Network Configuration: Gigabit Ethernet (\texttt{eth0}), LAN IP \texttt{11.12.5.8}.
        \item Virtual Interface: \texttt{tun0} (\texttt{10.8.0.1/24}, MTU 1400).
        \item Firewall Configuration: \texttt{iptables} NAT masquerade active with \texttt{net.ipv4.ip\_forward=1}.
    \end{itemize}
    \item \textbf{Packet Monitoring and Instrumentation Node}:
    \begin{itemize}
        \item Network Sniffer: Promiscuous packet capture executing Wireshark 4.2.2 and \texttt{tcpdump}.
        \item Taps: Dual-homed monitoring on both external Wi-Fi transport (TCP 7000) and internal server virtual interface (\texttt{tun0}).
    \end{itemize}
\end{itemize}

\section{Unit Testing Framework and Execution}
\label{sec:unit_testing_execution}

The project repository includes six automated unit and integration test suites located in \texttt{tests/}.

\subsection{Cryptographic Primitive Verification}
\begin{enumerate}
    \item \textbf{ML-KEM-1024 Verification (\texttt{tests/test\_mlkem.py})}: Validates keypair generation ($pk = 1,568$\,B, $sk = 3,168$\,B), encapsulation ($c = 1,568$\,B, $SS = 32$\,B), and decapsulation. Verifies that modifying any single bit of ciphertext $c$ produces an invalid shared secret, satisfying IND-CCA2 failure properties.
    \item \textbf{ML-DSA-65 Verification (\texttt{tests/test\_mldsa.py})}: Validates signing and verification over varying message sizes ($pk = 1,952$\,B, $\sigma = 3,309$\,B). Validates that tampering with either the message or the signature causes \texttt{verify()} to return \texttt{False}.
    \item \textbf{Classical AEAD Verification (\texttt{tests/test\_classical.py})}: Validates that HKDF-SHA256 expands 32-byte shared secrets into identical 256-bit AES keys across both runtimes. Verifies AES-256-GCM encryption and decryption across 10,000 randomized plaintext packets.
    \item \textbf{Dynamic IP Pool Verification (\texttt{tests/test\_ip\_pool.py})}: Tests concurrent allocation and deallocation across 253 virtual client leases (\texttt{10.8.0.2} to \texttt{10.8.0.254}), confirming zero address collisions or memory leaks under multi-threaded load.
\end{enumerate}

\begin{lstlisting}[basicstyle=\ttfamily\scriptsize, caption={Automated Test Suite Console Execution Output.}]
======================================================================
  QuantumSafeVPN -- Security & Integration Test Suite Execution
======================================================================
[TEST 1] tests.test_mlkem         ... [OK] (ML-KEM-1024 Key Encap/Decap)
[TEST 2] tests.test_mldsa         ... [OK] (ML-DSA-65 Signatures & Tamper Rejection)
[TEST 3] tests.test_classical     ... [OK] (HKDF-SHA256 & AES-256-GCM Engine)
[TEST 4] tests.test_vpn_tampering ... [OK] (Packet Tampering & Auth Tag Validation)
[TEST 5] tests.test_ip_pool       ... [OK] (Dynamic IP Pool Leases 10.8.0.2-254)
[TEST 6] tests.test_integration   ... [OK] (End-to-End PQC Handshake & Dispatch)
----------------------------------------------------------------------
Ran 6 test suites in 4.182s
STATUS: ALL TESTS PASSED (100% SUCCESS RATE)
======================================================================
\end{lstlisting}

\section{Wire Protocol Interoperability and Cross-Platform Testing}
\label{sec:wire_interoperability_testing}

To verify cross-platform wire compatibility between the Android Kotlin client and the Linux Python server, a dedicated synthetic harness was executed.

Table~\ref{tab:wire_interop_matrix} summarizes the observed wire parameters across each handshake phase, demonstrating 100\% mathematical and binary alignment.

\begin{table}[htbp]
  \centering
  \small
  \caption{Wire Protocol Interoperability and Alignment Verification Matrix.}
  \label{tab:wire_interop_matrix}
  \begin{tabular}{@{}P{3.0cm}P{4.2cm}P{4.2cm}P{2.6cm}@{}}
    \toprule
    \textbf{Protocol Stage} & \textbf{Linux Server (Python)} & \textbf{Android Client (Kotlin)} & \textbf{Status} \\
    \midrule
    \textbf{Length Framing} & 4-byte big-endian (\texttt{!I}) & 4-byte big-endian (\texttt{putInt}) & 100\% Identical \\
    \textbf{Stage 1: Server Hello} & Sends 1,952B ML-DSA-65 PK & Receives 1,952B; pins hash & 100\% Identical \\
    \textbf{Stage 2: Client Hello} & Receives 1,568B ML-KEM PK & Generates \& sends 1,568B PK & 100\% Identical \\
    \textbf{Stage 3A: Ciphertext} & Encapsulates; sends 1,568B $c$ & Decapsulates matching 32B $SS$ & 100\% Identical \\
    \textbf{Stage 3B: Signature} & Signs transcript (3,309B $\sigma$) & Verifies $\sigma$ over transcript & 100\% Identical \\
    \textbf{Stage 3C: Config JSON} & Sends \texttt{\{"assigned\_ip":"10.8.0.2"\}} & Parses JSON; configs TUN & 100\% Identical \\
    \textbf{Stage 4: Key KDF} & HKDF-SHA256 ($SS$, salt=$0^{32}$) & HKDF-SHA256 ($SS$, salt=$0^{32}$) & 100\% Identical \\
    \textbf{Data Framing} & \texttt{[4B len] + [12B IV] + Payload + [16B tag]} & \texttt{[4B len] + [12B IV] + Payload + [16B tag]} & 100\% Identical \\
    \bottomrule
  \end{tabular}
\end{table}

\section{Security and Adversary Injection Testing}
\label{sec:security_adversary_testing}

Security testing evaluated system resilience under active network attack conditions.

\subsection{Ciphertext Bit-Flipping and Tag Validation (\texttt{test\_vpn\_tampering.py})}
An adversary injection proxy was placed between the client and server. The proxy intercepted valid encrypted data frames traversing TCP port 7000 and introduced random bit-flips at varying offsets:
\begin{enumerate}
    \item \textbf{IV Tampering}: Flipping a single bit in the 12-byte initialization vector.
    \item \textbf{Ciphertext Tampering}: Flipping bits within the encrypted payload body.
    \item \textbf{Authentication Tag Tampering}: Modifying bytes within the 16-byte Galois MAC tag.
\end{enumerate}
In all 1,000 executed tampering trials, the receiving endpoint immediately raised an \texttt{AEADBadTagException} (on Android) or \texttt{InvalidTag} exception (on Linux), instantly dropping the frame without writing any corrupted bytes into the TUN device. Zero bit-flipped packets bypassed the authentication boundary.

\subsection{Anti-Replay Counter Injection Testing}
An adversary proxy captured 100 valid encrypted data frames (with sequence counters $S \in [50, 150]$) and re-injected them into the TCP stream 10 seconds later (after the active session counter had progressed to $S = 450$).
\begin{itemize}
    \item \textbf{Result}: The receiver's sequence check evaluated $S_{\text{incoming}} \le S_{\text{last\_received}}$ ($50 \le 450$), successfully dropping all 100 replayed packets with a security warning log.
\end{itemize}

\subsection{Man-in-the-Middle Handshake Tampering}
An active proxy substituted the server's ML-DSA-65 public key during Stage~1 with a rogue keypair generated by the attacker.
\begin{itemize}
    \item \textbf{Result}: The Android client evaluated the received key against the pinned SHA-256 fingerprint stored in \texttt{SecurityManager}. The fingerprint mismatch immediately triggered a \texttt{SecurityException}, aborting the connection prior to transmitting any client key encapsulation data.
\end{itemize}

\section{Full-Tunnel Routing Verification and DNS Leak Testing}
\label{sec:dns_leak_routing_testing}

To verify that user traffic cannot leak around the VPN tunnel:
\begin{enumerate}
    \item \textbf{Default Route Verification}: Executing \texttt{ip route show} and inspect network routing tables via ADB shell confirmed that the default routing table directed all traffic (\texttt{0.0.0.0/0}) into the virtual \texttt{tun0} interface.
    \item \textbf{External IP Verification}: While connected, querying public IP lookup services (\texttt{https://api.ipify.org}) returned the Linux gateway's public IP address rather than the mobile device's underlying Wi-Fi or cellular IP.
    \item \textbf{DNS Leak Audit}: The client configured recursive DNS servers to \texttt{1.1.1.1} and \texttt{8.8.8.8}. Continuous packet capture on the physical Wi-Fi interface confirmed that zero cleartext UDP port 53 datagrams egressed through \texttt{wlan0}; 100\% of DNS queries were encapsulated within the encrypted TCP tunnel.
\end{enumerate}

\section{Comprehensive Test Execution Matrix}
\label{sec:comprehensive_test_matrix}

Table~\ref{tab:test_execution_matrix} documents the comprehensive test execution catalogue, comprising 20 formal test cases across functional, cryptographic, network, and UI domains.

\begin{longtable}{@{}P{1.4cm}P{2.6cm}P{3.4cm}P{3.2cm}P{3.2cm}P{1.2cm}@{}}
  \caption{Comprehensive System Test Execution Matrix.} \label{tab:test_execution_matrix} \\
  \toprule
  \textbf{Test ID} & \textbf{Category} & \textbf{Test Objective} & \textbf{Procedure} & \textbf{Observed Result} & \textbf{Status} \\
  \midrule
  \endfirsthead
  \multicolumn{6}{c}{\small\textit{Table \ref{tab:test_execution_matrix} continued from previous page}} \\
  \toprule
  \textbf{Test ID} & \textbf{Category} & \textbf{Test Objective} & \textbf{Procedure} & \textbf{Observed Result} & \textbf{Status} \\
  \midrule
  \endhead
  \bottomrule
  \endfoot

  \textbf{TC-01} & Permissions & Verify \texttt{VpnService} OS consent & Launch app; tap Connect for first time & System modal displayed; consent granted & \textbf{PASS} \\
  \addlinespace
  \textbf{TC-02} & Socket & Verify socket protection & Inspect socket file descriptor via \texttt{protect()} & Socket tagged; egresses via physical adapter & \textbf{PASS} \\
  \addlinespace
  \textbf{TC-03} & Crypto & ML-KEM-1024 keypair generation & Execute \texttt{generateMlkemKeypair()} & Produced 1568B PK and 3168B SK & \textbf{PASS} \\
  \addlinespace
  \textbf{TC-04} & Crypto & ML-KEM-1024 shared secret parity & Encap on server, decap on Android & Both recover identical 32B shared secret & \textbf{PASS} \\
  \addlinespace
  \textbf{TC-05} & Crypto & ML-DSA-65 identity signing & Server signs transcript; client verifies & Signature (3309B) verified successfully & \textbf{PASS} \\
  \addlinespace
  \textbf{TC-06} & Crypto & Keccak SHA3-256 cross-runtime & Hash identical 64B seed in Kotlin and Python & Produced bit-for-bit identical 32B digests & \textbf{PASS} \\
  \addlinespace
  \textbf{TC-07} & Crypto & HKDF-SHA256 session key derivation & Derive key from shared secret on both nodes & Both endpoints derive identical 256-bit key & \textbf{PASS} \\
  \addlinespace
  \textbf{TC-08} & Crypto & AES-256-GCM authenticated enc/dec & Encrypt 1500B IP packet; decrypt & Ciphertext decrypted cleanly; tag matches & \textbf{PASS} \\
  \addlinespace
  \textbf{TC-09} & Security & Ciphertext bit-flip detection & Flip bit at byte offset 42 in frame & \texttt{AEADBadTagException}; frame dropped & \textbf{PASS} \\
  \addlinespace
  \textbf{TC-10} & Security & Anti-replay monotonic enforcement & Replay packet with seq 45 when current is 90 & Packet dropped; replay warning logged & \textbf{PASS} \\
  \addlinespace
  \textbf{TC-11} & Security & Server identity pinning check & Present rogue ML-DSA keypair & Fingerprint mismatch; connection aborted & \textbf{PASS} \\
  \addlinespace
  \textbf{TC-12} & Network & Dynamic IP pool lease allocation & Connect client; inspect assigned IP & Assigned \texttt{10.8.0.2}; subnet \texttt{255.255.255.0} & \textbf{PASS} \\
  \addlinespace
  \textbf{TC-13} & Network & Linux kernel NAT packet forwarding & Send ping to \texttt{1.1.1.1} from mobile & Packet masqueraded via \texttt{eth0}; reply OK & \textbf{PASS} \\
  \addlinespace
  \textbf{TC-14} & Network & MTU fragmentation boundary & Transmit 1400-byte ICMP payload & Packet forwarded without IP fragmentation & \textbf{PASS} \\
  \addlinespace
  \textbf{TC-15} & Network & DNS leak prevention & Query \texttt{dnsleaktest.com} & Zero ISP leaks; resolver shows \texttt{1.1.1.1} & \textbf{PASS} \\
  \addlinespace
  \textbf{TC-16} & UI & State badge transition sequence & Tap Connect; observe UI flow & DISCONNECTED $\to$ CONNECTING $\to$ CONNECTED & \textbf{PASS} \\
  \addlinespace
  \textbf{TC-17} & UI & Telemetry throughput updating & Stream 100\,MB file; observe UI counters & Uplink/downlink rates update in real time & \textbf{PASS} \\
  \addlinespace
  \textbf{TC-18} & Storage & Server configuration persistence & Update port to 7001; restart app & Settings retrieved from Android Keystore & \textbf{PASS} \\
  \addlinespace
  \textbf{TC-19} & Lifecycle & Graceful disconnection teardown & Tap Disconnect; check system state & TUN PFD released; socket closed cleanly & \textbf{PASS} \\
  \addlinespace
  \textbf{TC-20} & Resilience & Abrupt Wi-Fi disconnect handling & Disable Wi-Fi while tunnel active & Service handles reset; UI displays ERROR & \textbf{PASS} \\
\end{longtable}

All 20 test cases achieved a 100\% pass rate, confirming the functional validity, cryptographic integrity, and operational robustness of QuantumSafe VPN.
